
\documentclass[11pt]{article}
\usepackage[margin=1in]{geometry}
\usepackage{amsmath,amssymb,booktabs,hyperref,graphicx,parskip,enumitem}
\usepackage{xcolor}
\hypersetup{colorlinks=true,linkcolor=blue!50!black,urlcolor=blue!50!black}
\setlist{nosep}

\title{Replication Report: (title not recorded)\\[2pt]\large OSTI 2475938 \textbar{} Coverage 6/10, Agreement 7/10}
\author{Rick Stevens \and Ollie (AI Assistant)\\\small Argonne National Laboratory}
\date{2026-04-27}

\begin{document}\maketitle

\section{Source paper}
\begin{itemize}
\item \textbf{Title:} (title not recorded)
\item \textbf{Authors:} Not recorded
\item \textbf{Year:} n/a
\item \textbf{Domain:} Not recorded
\item \textbf{Identifier:} OSTI 2475938
\end{itemize}


\section{Replication objective}
The central claim of the paper concerns not recorded. Our objective was to reproduce the headline numerical/qualitative results within the constraints of the AI-driven Replication Project (24--72\,h, commodity GPU/CPU compute, no author contact). A replication plan is archived at \texttt{2475938-Updated-Virophage-Taxonomy-and-Distinction-from/replication\_plan.pdf}.

\section{Methodology}
We followed the standard Replication Project workflow: ingest the source PDF, extract method and data pointers, draft a replication plan, attempt the smallest end-to-end pipeline that produces a comparable number, then scale up where compute and data permit. Tools used in this replication are catalogued in the meta-paper's computational tools inventory (see \texttt{COMPUTATIONAL\_TOOLS\_INVENTORY.md}).

This paper went through the Tier-Lift pipeline; supplementary notes at \texttt{.openclaw/workspace/24h-progress/tier-lift-v2/2475938.md}.

\subsection*{Paper-directory README excerpt}
\small
\par\medskip\textbf{Updated Virophage Taxonomy and Distinction from Polinton-like Viruses}\par
\par$\bullet$\ **OSTI ID:** 2475938
\par$\bullet$\ **Paper:** Roux et al., *Biomolecules* 13:204 (2023)
\par$\bullet$\ **Rank:** \#9
\par$\bullet$\ **Original replication-score target:** 9/10
\par$\bullet$\ **Realised score (this run):** \textasciitilde{}6/10 (simplified reference set)

\par\medskip\textbf{\# Why This Paper}\par
This genomics paper uses publicly available genome sequences, computational analysis, and quantitative taxonomic criteria, allowing for fully automated AI replication and validation.

\par\medskip\textbf{\# Replication Plan (high level)}\par
Collect genome datasets → predict ORFs (Prodigal) → apply HMM profiles for marker genes (MCP, ATPase, PRO, penton + PLV HMM) → align \& trim markers → build ML phylogenies (IQ-TREE) → confirm virophage monophyly and virophage/PLV discrimination.

\par\medskip\textbf{\# Status}\par
\par$\bullet$\ [x] Paper reviewed
\par$\bullet$\ [x] Data identified (22 NCBI GenBank accessions; 13 retained after size filtering)
\par$\bullet$\ [x] Code implemented (Prodigal + HMMER + MAFFT + trimAl + IQ-TREE)
\par$\bullet$\ [x] Results validated (congruent 4-marker phylogeny; clean HMM separation of virophages vs PLV/adintovirus)

\par\medskip\textbf{\# This Replication}\par

Run on **CherryRd** (macOS, conda env `virophage`), total wall time \textless{} 20 min compute
(+ \textasciitilde{}5 min download). Deliverables:
\par$\bullet$\ **Report PDF:** [`replication/report/report.pdf`](replication/report/report.pdf) (6 pages)
\par$\bullet$\ **Figures:** `replication/report/phylogeny\_4markers.pdf`, `replication/report/hmm\_heatmap.pdf`
\par$\bullet$\ **Classification table:** `replication/results/hmm\_summary.tsv`
\par$\bullet$\ **Trees (Newick):** `replication/results/trees/\{MCP,ATPase,PRO,Penton\}.treefile`
\par$\bullet$\ **Alignments:** `replication/results/markers/\{MCP,ATPase,PRO,Penton\}.\{aln,trim\}`

\par\medskip\textbf{\#\# Key findings (replicated)}\par

1. **HMM classification matches the paper.** 7/13 genomes (Sputnik, Mavirus,
   SW01, YSLV1–4) hit all four virophage marker HMMs strongly (MCP ≥ 472 bits);
   the remaining 6 (TVSG\_01 PLV, three *Chrysochromulina parva* virophages,
   adintovirus, and two partial/mis-labelled GenBank records) have **zero**
   hits above inclusion to any of the 19 virophage HMMs or the PLV HMM.
2. **Virophage–PLV distinction is reproducible.** Virophages score
   ≥ 82–765 bits on the four morphogenesis markers; PLVs / adintoviruses
   score 0 — directly confirming the paper's statement that the marker
   HMMs separate the two classes.
3. **Four-marker phylogenetic congruence.** MCP, ATPase, PRO, and penton

[\textbackslash{}textit\{README truncated\}]
\normalsize

The replication was driven by an OpenClaw agent loop (Claude Opus / GPT-5 class models) issuing shell, Python, and (where applicable) GPU jobs. Compute usage and wall-time for individual papers are summarised in the meta-paper's resource table; not separately recorded for this paper unless noted in the Tier-Lift artefact. Sub-agents were spawned for replication-plan generation and (for papers in batches 1--4) for tier-lift extension runs.

\section{What was reproduced}
\begin{itemize}
\item Not recorded.
\end{itemize}

\section{What was missing or incomplete}
\begin{itemize}
\item Not recorded.
\end{itemize}
The dominant blocker categories for this paper, mapped onto the meta-paper's 9-category friction taxonomy (\S4), are listed in \S\ref{sec:friction} below.

\section{Quantitative agreement}
Per-quantity numerical comparisons are not separately tabulated in the structured evaluation record for this paper beyond the agreement rationale below; where the rationale references specific numbers, they are reproduced verbatim. A full numerical comparison table is recorded in the per-replication artefacts (see \S\ref{sec:repro}). For papers below 8/8, the agreement rationale identifies the specific quantities that diverge.

\paragraph{Agreement rationale.} Not recorded.

\section{Score and friction-point tagging}\label{sec:friction}
\begin{itemize}
\item \textbf{Coverage:} 6/10 --- Not recorded.
\item \textbf{Agreement:} 7/10 --- Not recorded.
\end{itemize}

\noindent\textbf{Friction points encountered (taxonomy tags):}
\begin{itemize}
\item \textbf{F5}: Force-field / hyperparameter drift (unspecified configuration)
\end{itemize}

\section{Follow-on questions}
\begin{itemize}
\item Not recorded.
\end{itemize}

\section{Reproducibility pointers}\label{sec:repro}
\begin{itemize}
\item \textbf{Paper directory:} \texttt{2475938-Updated-Virophage-Taxonomy-and-Distinction-from}
\item \textbf{Replication artefacts:}
\begin{itemize}
\item \texttt{/Users/stevens/Dropbox/REPLICATE-PROJECT/2475938-Updated-Virophage-Taxonomy-and-Distinction-from/replication}
\end{itemize}
\item \textbf{Replication-dir hint from JSONL:} \texttt{/Users/stevens/Dropbox/REPLICATE-PROJECT/2475938-Updated-Virophage-Taxonomy-and-Distinction-from}
\item \textbf{Status:} tier\_lift\_v2\_2026\_04\_27
\end{itemize}

To re-run, change directory into the paper directory above and consult its \texttt{README.md} for the entry-point script. Most replications follow the convention \texttt{replication/run\_*.sh} or \texttt{replication/<subdir>/run.py}.

\end{document}
